\subsection{\large Archivo Principal \texttt{(chat.c)}}
Este archivo es el orquestador de toda la lógica de la conexión P2P y la interfaz gráfica que se muestra al usuario. Al modular ambas partes, se requiere la bifurcación del sistema en dos procesos independientes (un padre y un hijo) mediante la directiva \texttt{fork()}, lo que permite manejar la aplicación completa de forma concurrente sin que la red bloquee la interfaz.

\begin{lstlisting}[style=estiloC]
#include <stdio.h>
#include <stdlib.h>
#include <unistd.h>
#include <sys/wait.h>
#include "header.h"

// Tubería que se usa para comunicar los dos procesos (Pipe A y Pipe B)
Tuberia principal;

int main(int argc, char *argv[])
{
    // Creamos las dos tuberías a utilizar durante la ejecución
    if (pipe(principal.B) < 0 || pipe(principal.A) < 0)
    {
        perror("Error al crear los pipes");
        return 1;
    }
    
    // Creamos un proceso hijo mediante fork para levantar la conexión P2P
    pid_t process = fork();
    
    // Si ocurre algún error en la bifurcación, se cancela la ejecución
    if (process < 0)
        return 1;
        
    // De lo contrario, el proceso hijo (process == 0) se encarga de la conexión P2P
    if (process == 0)
    {
        // El hijo cierra los extremos no usados de las tuberías
        close(principal.A[0]);
        close(principal.B[1]);
        
        // Llama a la función que levanta la red y los sockets
        conexion(&principal);
        
        // Cierra los extremos utilizados para terminar exitosamente y liberar memoria
        close(principal.A[1]); 
        close(principal.B[0]);

        exit(0);
    }
    else // El proceso padre se encarga de renderizar la interfaz gráfica
    {
        // Se duerme 2 segundos para dar tiempo a que el back-end levante el servidor
        sleep(2);
        
        // El padre cierra los extremos que no utiliza de las tuberías
        close(principal.A[1]);
        close(principal.B[0]);
        
        // Llama a la función principal de GTK que crea la vista
        interfaz_grafica(argc, argv, &principal);
        
        // Cierra los extremos utilizados tras cerrar la ventana
        close(principal.A[0]); 
        close(principal.B[1]);
        
        // El padre espera a que el proceso hijo termine para cerrarse exitosamente
        wait(NULL);
    }
    return 0;
}
\end{lstlisting}